\section{Preliminaries}
\label{sec:preliminaries}

Let $\T =
\R/\Z$ denote the cycle, i.e., the additive group of reals modulo
$1$. We also denote by $\T_q$ its cyclic subgroup of order $q$, i.e.,
the subgroup given by $\{0,1/q,\ldots,(q-1)/q\}$.

For two probability distributions $P,Q$ over some discrete domain, we
define their statistical distance as $\sum |P(i)-Q(i)|/2$ where~$i$
ranges over the distribution domain, and extend this to continuous
distributions in the obvious way. We recall the following easy fact
(see, e.g.,~\cite[Eq.~(2.3)]{AldousDiaconis87} for a proof).
\begin{claim}\label{clm:separationdist}
  If $P$ and $Q$ are two probability distributions such that~$P(i) \ge
  (1-\eps)Q(i)$ holds for all~$i$, then the statistical distance
  between $P$ and $Q$ is at most~$\eps$.
\end{claim}



We will use the following immediate corollary of the leftover hash
lemma~\cite{DBLP:journals/siamcomp/HastadILL99}.

\begin{lemma}
  \label{lem:lhl}
  Let $k, n, q \ge 1$ be integers, and $\epsilon > 0$ be such that $n
  \ge k \log_2 q + 2 \log_2 (1/\epsilon)$. For $\mx{H} \gets \T_q^{k \times
    n}$, $\vc{z} \gets \binset^n$, $\vc{u} \gets \bbT_q^k$, the
  distributions of $(\mx{H}, \mx{H}\vc{z})$ and $(\mx{H}, \vc{u})$ are
  within statistical distance at most $\epsilon$.
\end{lemma}


A \emph{distinguishing problem} $P$ is defined by two distributions $P_0$ and $P_1$, and a solution to the problem is the ability to distinguish between these distributions. The \emph{advantage} of an algorithm $\cA$ with binary output on $P$ is defined as
\[
\adv[\cA] = \abs{\Pr[\cA(P_0)] - \Pr[\cA(P_1)]}~.
\]
A reduction from a problem~$P$ to a problem~$Q$ is an efficient (i.e.,
polynomial-time) algorithm~$\cA^{\cB}$ that solves~$P$ given
access to an oracle~$\cB$ that solves~$Q$. Most of our reductions
(in fact all except the one in Lemma~\ref{lem:lelwe}) are what we 
call ``transformation reductions:'' these reductions perform some transformation
to the input and then apply the oracle to the result. 


\subsection{Lattices}
\label{sec:lattices}

An $n$-dimensional (full-rank) lattice $\Lambda \subseteq \R^n$ is the
set of all integer linear combinations of some set of $n$ linearly
independent \emph{basis} vectors $\matB = \set{\vecb_{1}, \ldots,
  \vecb_{n}} \subseteq \R^n$,
\[ \Lambda = \lat(\matB) = \set[\Big]{
  \sum_{i \in [n]} z_i \vecb_i\; :\; \vecz \in \mathbb{Z}^n }. \]
The \emph{dual lattice} of $\Lambda \subset \R^n$ is defined as
$\Lambda^* = \set{ \vecx \in \R^n : \inner{\Lambda,\vecx} \subseteq \Z}$.

The \emph{minimum distance} (or \emph{first successive minimum}) $\lambda_1(\Lambda)$ of a lattice
$\Lambda$ is the length of a shortest nonzero lattice vector, i.e., $\lambda_1(\Lambda) = \min_{\veczero \neq
  \vecx \in \Lambda} \length{\vecx}$.  For an approximation ratio $\gamma = \gamma(n) \ge 1$, the $\gapsvp_\gamma$ is the problem of deciding,
	given a basis $\matB$ of an $n$-dimensional lattice $\Lambda = \lat(\matB)$ and a number $d$, between the case where $\lambda_1(\lat(\matB)) \le d$ 
	and the case where $\lambda_1(\lat(\matB)) > \gamma d$.
We refer to~\cite{KhotLLL25,RegevLLL25} for a recent account 
on the computational complexity of $\gapsvp_\gamma$. \dnote{I added that sentence}


\subsection{Gaussian measures}
\label{sec:gaussian_measures}


For $r > 0$, the $n$-dimensional Gaussian function $\rho_r : \R^{n} \to (0,1]$ is
defined as
\[ \rho_r(\vecx) := \exp(-\pi
\length{\vecx}^{2}/r^2).
\]
We extend this definition to sets, i.e., $\rho_r(A) = \sum_{\vecx \in A} \rho_r(\vecx) 
\in [0,+\infty]$ for any~$A \subseteq \R^{n}$.
The (spherical) continuous Gaussian distribution~$D_r$ is the distribution with density
function proportional to $\rho_r$. More generally, for a matrix $\matB$,
we denote by $D_\matB$ the distribution of $\matB \vecx$ where $\vecx$ is sampled
from $D_1$. When $\matB$ is nonsingular, its probability density function is
proportional to
\[
\exp(-\pi \vecx^T (\matB \matB^T)^{-1} \vecx).
\]
A basic fact is that for any matrices $\matB_1,\matB_2$, the sum of a sample from $D_{\matB_1}$ and an independent sample
from~$D_{\matB_2}$ is distributed like $D_{\matC}$ for $\matC = (\matB_1^{} \matB_1^T + \matB_2^{} \matB_2^T)^{1/2}$.

For an $n$-dimensional lattice $\Lambda$ and a vector $\vecu \in \R^n$, we define the \emph{discrete Gaussian distribution}~$D_{\Lambda+\vecu,r}$
as the discrete distribution with support on the coset $\Lambda+\vecu$ whose probability mass function is proportional to
$\rho_r$. There exists an efficient procedure that samples within negligible statistical distance of any (not too narrow) discrete Gaussian distribution 
(\cite[Theorem~4.1]{DBLP:conf/stoc/GentryPV08}; see also~\cite{DBLP:conf/crypto/Peikert10}). In the next lemma, proved in Section~\ref{sec:exactgpv},
we modify this sampler so that the output is distributed exactly as a discrete Gaussian. This also allows us to sample from
slightly narrower Gaussians. Strictly speaking, the lemma is not needed for our results,
and we could use instead the original sampler from \cite{DBLP:conf/stoc/GentryPV08}. Using our exact sampler leads to slightly
cleaner proofs as well as a (miniscule) improvement in the parameters of our reductions, and we include it here mainly in the hope that it finds
further applications in the future.

\begin{lemma}\label{lem:gpv}
There is a probabilistic polynomial-time algorithm that, given a basis $\matB$ of an $n$-dimensional
lattice $\Lambda = \lat(\matB)$, $\vecc \in \R^n$, and a parameter $r \ge \|\widetilde \matB\| \cdot \sqrt{\ln(2n+4)/\pi}$,
outputs a sample distributed according to $D_{\Lambda+\vecc,r}$. 
\end{lemma}

Here, $\widetilde \matB$ denotes the Gram-Schmidt orthogonalization of $\matB$,
and $\|\widetilde \matB\|$ is the length of the longest vector in it.
We recall the definition of the \emph{smoothing parameter} from~\cite{DBLP:journals/siamcomp/MicciancioR07}.

\begin{definition}
  \label{def:smoothing}
  For a lattice $\Lambda$ and positive real $\epsilon > 0$, the
  smoothing parameter $\smootheps(\Lambda)$ is the smallest real $s >
  0$ such that $\rho_{1/s}(\Lambda^* \setminus \set{\veczero}) \leq
  \epsilon$.
\end{definition}

\begin{lemma}[{{\cite[Lemma~3.1]{DBLP:conf/stoc/GentryPV08}}}]\label{lem:boundonsmoothing}
For any $\eps>0$ and $n$-dimensional lattice $\Lambda$ with basis $\matB$,
\[
\smootheps(\Lambda) \le \|\widetilde \matB \| \sqrt{\ln(2n(1+1/\eps))/\pi}.
\]
\end{lemma}


We now collect some known facts on Gaussian distributions and lattices.

\begin{lemma}[{{\cite[Lemma 4.1]{DBLP:journals/siamcomp/MicciancioR07}}}]
  \label{lem:smoothingcontinuous}
  For any $n$-dimensional lattice $\Lambda$, $\epsilon>0$, $r \geq
  \eta_{\epsilon}(\Lambda)$, the distribution of $\vecx \bmod \Lambda$ where $\vecx \leftarrow D_{r}$ is
	within statistical distance $\eps/2$ of the uniform distribution on cosets of $\Lambda$.
\end{lemma}

\begin{lemma}[{{\cite[Claim 3.8]{DBLP:journals/jacm/Regev09}}}]
  \label{lem:smoothing}
  For any $n$-dimensional lattice $\Lambda$, $\epsilon>0$, $r \geq
  \eta_{\epsilon}(\Lambda)$, and $\vecc \in \R^n$, we have $\rho_{r}(\Lambda + \vecc) \in [
  \tfrac{1-\epsilon}{1+\epsilon}, 1 ] \cdot \rho_{r}(\Lambda)$.
\end{lemma}

\begin{lemma}[{{\cite[Claim~3.9]{DBLP:journals/jacm/Regev09}}}]
\label{le:discrete_plus_cont_Gauss}
Let~$\Lambda$ be an $n$-dimensional lattice, let~$\vec{u}\in \R^n$ be
arbitrary, let $r, s > 0$ and let $t=\sqrt{r^2+ s^2}$. Assume that
$rs/t=1/\sqrt{1/r^2+1/s^2} \geq \eta_{\eps} (\Lambda)$ for some~$\eps <
1/2$. Consider the continuous distribution~$Y$ on~$\R^n$ obtained by
sampling from~$D_{\Lambda+\vec{u},r}$ and then adding a noise vector taken
from~$D_s$. Then, the statistical distance between~$Y$ and~$D_t$ is at
most~$4\eps$.
\end{lemma}

\begin{lemma}[{{\cite[Corollary~3.10]{DBLP:journals/jacm/Regev09}}}]
\label{lem:smoothip}
Let~$\Lambda$ be an $n$-dimensional lattice, let~$\vc{u}, \vc{z}\in \R^n$ be arbitrary, and let $r, \alpha > 0$. Assume that
$\parens{1/r^2+\left(\|{\vc{z}}\|/\alpha\right)^2}^{-1/2} \geq \eta_{\eps} (\Lambda)$ for some~$\eps <
1/2$. Then the distribution of $\langle\vc{z},\vc{v}\rangle+e$ where $\vc{v} \getsr D_{\Lambda+\vc{u}, r}$ and $e \getsr D_{\alpha}$, is within statistical distance $4\eps$ of $D_{\beta}$ for $\beta = \sqrt{(r \|\vc{z}\|)^2+\alpha^2}$.
\end{lemma}

\begin{lemma}[{{Special case of~\cite[Theorem 3.1]{DBLP:conf/crypto/Peikert10}}}]
\label{lem:peikertdiscretegauss}
Let $\Lambda$ be a lattice and $r,s>0$ be such that $s \ge \eta_\eps(\Lambda)$ for some $\eps \le 1/2$.
Then if we choose $\vecx$ from the continuous Gaussian $D_{r}$ and then choose $\vecy$ from the
discrete Gaussian $D_{\Lambda-\vecx,s}$ then $\vecx+\vecy$ is within statistical distance $8 \eps$ of
the discrete Gaussian $D_{\Lambda,(r^2+s^2)^{1/2}}$.
\end{lemma}


\subsection{Learning with Errors}
\label{sec:computational_problems}

For integers $n,q \ge 1$, an integer vector $\vec{s}
\in \Z^n$, and a probability distribution $\phi$ on~$\R$, let
$A_{q,\vecs, \phi}$ be the distribution over $\T_q^n \times \T$
obtained by choosing $\veca \in \T_q^n$ uniformly at random and an
error term $e$ from $\phi$, and outputting the pair $(\veca, b =
\inner{\veca, \vecs}+e) \in \T_q^n \times \T$.

\begin{definition}
  \label{def:lwe}
  For integers $n,q \ge 1$, an error distribution $\phi$
  over $\R$, and a distribution $\calD$ over $\Z^n$, the (average-case) decision variant of the $\lwe$
  problem, denoted $\lwe_{n,q,\phi}(\calD)$, is to distinguish 
	given arbitrarily many independent samples,
  the uniform distribution over $\T_q^n \times \T$ from
  $A_{q,\vec{s},\phi}$ for a fixed $\vec{s}$ sampled from $\calD$. The variant where the algorithm only gets a bounded number of samples $m\in\bbN$ is denoted $\lwe_{n,m,q,\phi}(\calD)$.
\end{definition}

Notice that the distribution $A_{q,\vec{s},\phi}$ only depends on $\vec{s} \bmod q$, and so one can assume
without loss of generality that $\vec{s} \in \{0,\ldots,q-1\}^n$. Moreover, using a standard
random self-reduction, for any distribution over secrets $\calD$, one can reduce $\lwe_{n,q,\phi}(\calD)$ to
$\lwe_{n,q,\phi}(U(\{0,\ldots,q-1\}^n))$, and we will occasionally use $\lwe_{n,q,\phi}$ to denote the latter
(as is common in previous work). 
When the noise is a Gaussian with parameter $\alpha>0$, i.e., $\phi=D_\alpha$,
we use the shorthand $\lwe_{n,q,\alpha}(\calD)$.
Since the case when $\calD$ is uniform over $\{0,1\}^n$ plays an important
role in this paper, we will denote it by $\zolwe_{n,q,\phi}$ (and by~$\zolwe_{n,m,q,\phi}$ when the algorithm only gets~$m$ samples). Finally, as we show in the following lemma, 
one can efficiently reduce $\lwe$ to the case in which the secret
is distributed according to the (discretized) error distribution and is hence somewhat short. This latter
form of \lwe, known as the ``normal form,'' was first shown hard in~\cite{DBLP:conf/crypto/ApplebaumCPS09}
for the case of prime $q$.
Here we observe that the proof extends to non-prime $q$, the new technical ingredient being Claim~\ref{clm:invertiblemodq} below.

\begin{lemma}\label{lem:normalform}
  For any $q \ge 25$, $n,m \geq 1$, $\alpha>0$, $\eps < 1/2$ and~$s \geq
  \sqrt{\ln(2n(1+1/\epsilon)/\pi)}/q$, there is an efficient
  (transformation) reduction from~$\lwe_{n,m,q,\alpha}$ to
  $\lwe_{n,m',q,\alpha}(\calD)$
  where~$m' = m - (16n + 4\ln \ln q)$ and $\calD=D_{\Z^n,q(\alpha^2+s^2)^{1/2}}$, that 
	turns advantage $\zeta$ into an advantage of at least $(\zeta-8\eps)/4$.
	In particular, assuming $\alpha \geq \sqrt{\ln(2n(1+1/\epsilon)/\pi)}/q$, we can take $s=\alpha$,
	in which case $\calD=D_{\Z^n,\sqrt{2}q\alpha}$.
\end{lemma}

\begin{proof}
  Consider the first $16n + 4\ln\ln q$ samples $(\vec{a},b)$.
	Using Claim~\ref{clm:invertiblemodq}, with probability at least
	$1-2e^{-1} \ge 1/4$, we can efficiently find a subsequence of the samples
	such that the matrix $A_0 \in \Z_q^{n \times n}$ whose columns
	are formed by the $\vec{a}$ in the subset (scaled up by $q$)
	has an inverse $A_0^{-1} \in \Z_q^{n \times n}$ modulo $q$. If we cannot find such a subsequence, we abort. 
	Let $\vec{b}_0 \in \T^n$ be the vector formed by the corresponding $b$ in the subsequence.
	Let also $\vec{b}'_0 \in \T_q^n$ be $\vec{b}_0 + \vec{x}$ where $\vec{x}$ is
	chosen from $D_{q^{-1}\Z^n - \vec{b}_0,s}$. 
	(Notice that the coset $q^{-1}\Z^n - \vecb_0$ is well defined
because $\vecb_0$ is a coset of $\Z^{n} \subseteq q^{-1}\Z^n$.)
  From each of the remaining $m'$ samples $(\vec{a},b) \in \T_q^n \times \T$ we produce a pair
	\[
	  \big(\vec{a}' = A_0^{-1} \vec{a}, b' = b - \inner{A_0^{-1} \cdot q \vec{a}, \vec{b}'_0}\big) \in \T_q^n \times \T.
	\]
	We then apply the given $\lwe$ oracle to the resulting $m'$ pairs and output its result. 
		
	We now analyze the reduction. First notice that the construction of $A_0$ depends
	only on the $\vec{a}$ component of the input samples, and hence the probability of 
	finding it is the same in case the input is uniform and in case it consists of $\lwe$ samples.
	It therefore suffices in the following to show that there is a distinguishing gap conditioned
	on successfully finding an $A_0$. 
	To that end, first observe that if the input samples $(\veca,b)$ are uniform in $\T_q^n \times \T$
	then so are the output samples $(\veca',b')$.
	Next consider the case that the input samples are distributed according to $A_{q,\vec{s},D_{\alpha}}$
	for some $\vec{s} \in \Z^n$. Then since~$s \geq \eta_{\epsilon}(q^{-1}\Z)$ by Lemma~\ref{lem:boundonsmoothing},
	using Lemma~\ref{lem:peikertdiscretegauss} we get that $\vec{b}_0' = q^{-1} A_0^T \vec{s} + \vece_0$
	where $\vece_0$ is distributed within statistical distance $8\eps$ from $D_{q^{-1}\Z^n, (\alpha^2+s^2)^{1/2}}$.
	Therefore, for each output sample $(\veca',b')$ we have
	\[
	   b' = b - \inner{A_0^{-1} \cdot q \vec{a}, \vec{b}'_0} 
		    = \inner{\veca,\vecs}+ e - \inner{\veca, \vecs} - \inner{A_0^{-1} q \veca, \vece_0}
				= \inner{-q \vec{e}_0, \veca'}+e,
	\]
	where $e$ is an independent error from $D_\alpha$. Therefore, the output samples
	are distributed according to $A_{q,-q\vec{e}_0,D_{\alpha}}$, completing the proof.
\end{proof}


\begin{claim}\label{clm:invertiblemodq}
For any $q \ge 25$, $n \ge 1$, and $t_1 \ge 4, t_2 \ge 1$, given a sequence of $t_1 n+t_2 \ln\ln q$
vectors $\vec{a}_1,\vec{a}_2,\ldots$ chosen uniformly and independently from $\Z_q^n$,
except with probability $e^{-t_1 n /16} + e^{-t_2/4}$, there exists a subsequence 
of $n$ vectors such that the $n \times n$ matrix they form is invertible modulo $q$. 
Moreover, such a subsequence can be found efficiently. 
\end{claim}
\begin{proof}
We consider the following procedure. 
Let $k$ be a counter, initialized to $0$, indicating the number of vectors currently 
in the subsequence, and let $A \in \Z_q^{n \times k}$ be the matrix whose columns are formed by the current
subsequence. We also maintain a unimodular matrix~$U \in \Z^{n \times n}$, 
initially set to the identity, satisfying the invariant that~$U \cdot A \in \Z_q^{n \times k}$ 
has the following form:
  its top $k \times k$ submatrix is upper triangular with each
  diagonal coefficient coprime with~$q$; its bottom $(n-k) \times
  k$ submatrix is zero.
The procedure considers the vectors $\vec{a}_i$ one by one.
For each vector $\vec{a}$, if it is such that the gcd of the last~$n-k$ entries of~$U
  \vec{a}$, call it $g$, is coprime with~$q$, then it does the following:
	it adds $\vec{a}$ to the subsequence, 
	computes (using, say, the extended GCD algorithm) a unimodular matrix $V$ that acts as identity on the first $k$ coordinates 
	and for which the last $n-k$ coordinates of $VU\vec{a}$ are $(g,0,\ldots,0)$,
	replaces $U$ with $VU$, and increments $k$. 

It is easy to see that the procedure's output is correct if it reaches $k=n$.
It therefore suffices to analyze the probability that this event happens. 
For this we use the following two facts to handle the cases~$k<n-1$ and~$k = n-1$, respectively. First, the probability that
the gcd of two uniformly random numbers modulo $q$ is coprime with $q$ is 
\[
\prod_{p|q,~p~\text{prime}} (1-p^{-2}) \ge \prod_{p~\text{prime}} (1-p^{-2}) = \zeta(2)^{-1} \approx 0.61,
\]
where $\zeta$ is the Riemann zeta function. 
Second, the probability that one uniformly random number modulo $q$ is coprime with $q$
is $\varphi(q)/q$, where $\varphi$ is Euler's totient function. By~\cite[Theorem 8.8.7]{BachShallit}, 
this probability is at least $(e^\gamma \ln \ln q + 3/(\ln\ln q))^{-1}$ where $\gamma$ is Euler's constant, which for $q \ge 25$
is at least $(4 \ln \ln q)^{-1}$. 

Using the (multiplicative) Chernoff bound, the first fact, and the fact that $U\vec{a}$ is uniform in $\Z_q^n$ since $U$ is unimodular, we see that the probability that $k < n-1$ after considering $t_1 n$ vector is at most $e^{-t_1 n /16}$. Moreover, once $k=n-1$, using the second fact we get that the probability that after considering $t_2 \ln \ln q$ additional vectors we still have $k=n-1$ is at most $e^{-t_2/4}$.
\end{proof}



\myparagraph{Unknown (Bounded) Noise Rate}
We also consider a variant of $\lwe$ in which the amount of noise is some unknown $\beta \le \alpha$ (as opposed
to exactly $\alpha$), with $\beta$ possibly depending on the secret $\vecs$. As the following lemma shows, this 
does not make the problem significantly harder.

\begin{definition}\label{def:lelwe}
For integers $n,q \ge 1$ and $\alpha \in (0,1)$, $\lwe_{n,q,\le\alpha}$ is the problem of solving $\lwe_{n,q,\beta}$ for any $\beta = \beta(\vc{s}) \le \alpha$.
\end{definition}

\begin{lemma}\label{lem:lelwe}
Let $\cA$ be an algorithm for $\lwe_{n, m, q, \alpha}$ with advantage at least $\epsilon>0$. Then there exists an algorithm $\cB$ for $\lwe_{n,m',q,\le\alpha}$ using oracle access to~$\cA$ and with advantage at least $1/3$, where both $m'$ and its running time are $\poly(m,1/\epsilon, n, \log q)$.
\end{lemma}

The proof is standard (see, e.g., \cite[Lemma 3.7]{DBLP:journals/jacm/Regev09} for
the analogous statement for the search version of $\lwe$). The idea is to use Chernoff bound to estimate $\cA$'s success probability on the uniform distribution, and then add noise in small increments to our given distribution and estimate $\cA$'s behavior on the resulting distributions. If there is a gap between any of these and the uniform behavior, the input distribution is deemed non-uniform. The full proof is omitted.

\znote{Commented out proof is down here. Fix or remove.
%
%\begin{proof}
%\znote{Only sketch for now. Also I am not sure all the specific constants are correct, need to double check.}
%Let $\cA$ be as in the lemma statement. Note that $\adv[\cA] \ge \epsilon$ means that
%\begin{equation}
%\Pr_{\vc{s}}\left[\abs{\Pr[\cA(\mx{A}, (\mx{A}\cdot\vc{s})/q+\vc{e})] -  \Pr[\cA(\mx{A}, \vc{u})] \ge \epsilon/2 }\right] \ge \epsilon/2~.
%\end{equation}
%Namely, for $\epsilon/2$ fraction of the secret vectors we get a good distinguishing gap.
%
%We define the algorithm $\cB$ as follows: $\cB$ will draw many samples $(\mx{A}, \vc{d})$ from its input distribution (recall that either all are of the form $A_{\vc{s}, \beta}$ or all are uniform). It will add to each of these samples an additional noise from $D_{\gamma}$, for all $\gamma$ over a fine grid of polynomially many values between $0$ and $\alpha$.
%
%Then $\cB$ will estimate $\cA$'s acceptance probability on the uniform distribution and on each of the $\gamma$ distributions. The estimate is such that with probability $1-\epsilon/6$, no estimate deviates by more than $\epsilon/20$ from the correct value. This can be done in $\poly(1/\epsilon, n, T)$ time.
%
%Then, if either of the $\gamma$'s yielded a value that is different from the value of the uniform distribution by more than $\epsilon/8$, then output $1$. Otherwise output $0$.
%
%The analysis proceeds as follows: First, we note that when $\cB$ is given the uniform distribution, it will output $0$ with all but $\epsilon/6$ probability. Next, we analyze $\cB$'s behavior on $A_{\vc{s}, \beta}$.
%
%% consider the distributions $(\mx{A}, \mx{A}^T\vc{s}/q+\vc{e})$ for $\vc{e} \getsr D_{\sqrt{\beta^2+\gamma^2}}$ for a fine grid of $\gamma$ between $0$ and $\alpha$. Note that these distributions are easy to generate without knowing $\beta$.
%
%If we choose the grid of $\gamma$'s fine enough, then there exists $\gamma$ such that $D^m_{(\beta^2+\gamma^2)^{1/2}}$ and $D^m_{\alpha}$ are within $\epsilon/4$ statistical distance. This means that for this specific $\gamma$, the algorithm $\cA$ will have a distinguishing gap $\epsilon/4$ for at least an $\epsilon/2$ fraction of the $\vc{s}$. \znote{Is the statistical distance between two Gaussians bounded by the difference of their standard deviations or the difference of variances?}\onote{use Claim 2.2 of~\cite{DBLP:journals/jacm/Regev09} which bounds the statistical distance between two normal variables by $9(\beta/\alpha - 1)$ assuming $\alpha < \beta \le 2 \alpha$} This means that $\cB$ will output $1$ with probability at least $\epsilon/2 - \epsilon/6 \ge \epsilon/3$.
%
%This implies a distinguishing gap of at least $\epsilon/6$, which can be improved to $\epsilon/3$ if we notice that we double counted the event of erroneous estimates of $\cA$'s acceptance probability.
%\end{proof}
}






\myparagraph{Relation to Lattice Problems}
Regev~\cite{DBLP:journals/jacm/Regev09} and Peikert~\cite{DBLP:conf/stoc/Peikert09} showed quantum and classical reductions (respectively) from the worst-case hardness of the $\gapsvp$ problem to the \onote{added:}search version of $\lwe$. 
(We note that the quantum reduction in~\cite{DBLP:journals/jacm/Regev09} also shows a reduction from $\sivp$.)
As mentioned in the introduction, the classical reduction only works when the modulus $q$ is exponential
in the dimension~$n$.
This is summarized in the following theorem, which is derived from
{{\cite[Theorem~3.1]{DBLP:journals/jacm/Regev09}}} and
{{\cite[Theorem~3.1]{DBLP:conf/stoc/Peikert09}}}.

\begin{theorem}\label{thm:worstcaseaveragecase}
Let $n, q \ge 1$ be integers and let $\alpha \in (0,1)$ be such that $\alpha q \ge 2 \sqrt{n}$. Then there exists a quantum reduction from worst-case $n$-dimensional $\gapsvp_{\widetilde{O}(n/\alpha)}$ to $\lwe_{n,q,\alpha}$. If in addition $q \ge 2^{n/2}$ then there is also a classical reduction between those problems. 
\end{theorem}

In order to obtain hardness of the \emph{decision} version of \lwe, which is the one we consider throughout the paper, one employs a search-to-decision reduction. Several such reductions appear in the literature (e.g.,~\cite{DBLP:journals/jacm/Regev09,DBLP:conf/stoc/Peikert09,DBLP:conf/eurocrypt/MicciancioP12}). The most recent reduction by Micciancio and Peikert~\cite{DBLP:conf/eurocrypt/MicciancioP12}, which essentially subsumes all previous reductions, 
requires the modulus~$q$ to be smooth. Below we give the special case when the modulus is a power of~$2$, which suffices for our purposes. It follows from our results that (decision) $\lwe$ is hard not just for a smooth modulus~$q$, as follows from~\cite{DBLP:conf/eurocrypt/MicciancioP12}, but actually for all moduli~$q$, including prime moduli, with only a small deterioration in the noise
(see Corollary~\ref{cor:mod-reduction-2}).

\begin{theorem}[{{Special case of~\cite[Theorem 3.1]{DBLP:conf/eurocrypt/MicciancioP12}}}]\label{thm:searchtodecisionmicpei}
Let $q$ be a power of $2$, and $\alpha$ satisfy $1/q < \alpha < 1/\omega(\sqrt{\log n})$. Then there exists an efficient reduction from search $\lwe_{n,q,\alpha}$
to (decision) $\lwe_{n,q,\alpha'}$ for $\alpha'=\alpha \cdot \omega(\log n)$.
\end{theorem}





%%% Local Variables:
%%% mode: latex
%%% TeX-master: "lwehard"
%%% End:
